\section{Q4: How can research experience improve the researcher?}
\label{sec:learning}

A research system can improve its current solution without becoming better at conducting the next investigation. Q4 concerns this distinction: which information survives an investigation, what component it changes, and whether the change benefits subsequent decisions. We separate \emph{within-task adaptation}, reuse across runs on related tasks, and transfer to held-out task families. These are evaluation settings, not interchangeable descriptions of learning. A persistent database can support the second without demonstrating the third; updating model weights during one problem can still be entirely task-local. Research taste requires an additional test of which questions deserve attention, beyond solving assigned problems more efficiently.

\subsection{External memory, reusable skills, and local reflection}

AgentRxiv makes research reports available to later laboratories, while Agora represents shared hypotheses, results, and verifications through Git-based records~\cite{agentrxiv,agora}. Both treat research as an accumulating collective process. Their distinction from an ordinary conversation history is persistence and reuse across sessions. However, persistence alone does not identify the causal contribution of memory: later agents may receive better candidate solutions, more cumulative computation, or more informative task-specific evidence. Agora explicitly leaves discovery gains per unit of computation to matched comparisons.

EvoScientist separates ideation memory from experimentation memory, distilling prior trajectories into directions and implementation strategies. InternAgent-1.5 combines procedural strategies, within-trajectory episodes, and semantic knowledge accumulated across sessions~\cite{evoscientist,internagent-1-5}. These architectures suggest that a research memory should distinguish \emph{what happened}, \emph{why an action was selected}, and \emph{when a lesson applies}. This distinction is our design interpretation, rather than a demonstrated universal memory schema. AutoResearchClaw emphasizes lessons extracted from previous failures; its cross-run ablation primarily supports avoiding known failure modes rather than raising the scientific-quality ceiling~\cite{autoresearchclaw}.

Reusable skills go beyond storing a report by prescribing actions under recognizable conditions. HypoForge refines hypothesis-generation and testing skills using comparative critique and execution feedback, then evaluates them on held-out tasks without updating the backbone~\cite{hypoforge}. This is a useful form of strategy adaptation, but its test-pass metric includes significance for reference hypotheses; it does not independently establish false-positive control or question importance. A skill can learn how to satisfy an evaluator while remaining a poor guide to informative science.

Local feedback must also retain its own category. Reflexion stores verbal reflections without weight updates; Dolphin feeds experimental analysis into subsequent iterations; OR-Agent uses hierarchical reflections and backtracking during heuristic search~\cite{reflexion,dolphin,or-agent}. These methods motivate feedback-sensitive researchers, but their respective evidence should not be rewritten as general cross-task scientific learning. In particular, OR-Agent's verbal gradient and momentum terminology describes an analogy for reflection, not gradient updates to the language model.

\subsection{Changing the research program or training the model}

Code evolution changes the procedure that produces future actions. ADAS searches agent programs using an archive of designs and evaluations; Darwin G\"odel Machine maintains evolving coding agents that modify their own implementations~\cite{adas,darwin-godel-machine}. They are foundations for research-agent design, with evidence from specified benchmark tasks rather than scientific agenda selection. Bilevel Autoresearch makes the inner research loop itself the optimization target: an outer process analyzes traces and generates mechanisms that modify search~\cite{bilevel-autoresearch}. Its reported study concerns one small GPT-pretraining setting. It demonstrates a bounded outer-to-inner intervention, not unbounded recursive improvement or general transfer between research domains.

Parameter training changes the model that proposes or evaluates decisions. CycleResearcher uses an automated reviewer for iterative preference training of research generation~\cite{cycleresearcher}. Its core evaluation includes model-generated experimental results that were not executed; an appendix separately illustrates integration with execution. Consequently, improved manuscript or reviewer scores must not be interpreted as improved empirical discovery. The distinction matters for training data: a fluent retrospective explanation is not necessarily a record of an experiment that happened.

ResearchTrails extracts decision sequences from human Git histories and studies demonstrations, distilled skills, supervised learning, and reinforcement learning~\cite{researchtrails}. Its principal prediction target is the recorded next research decision. Commit histories provide incomplete evidence about human exploration, and similarity to a historical decision does not establish that an alternative policy would discover more. SciSense instead synthesizes citation-conditioned sensemaking trajectories, including reconstructions of paths toward known contributions, and distills them into planning models~\cite{scisense}. Such reconstructions offer structured supervision, but they are not observations of the authors' actual deliberations. Its downstream execution study and judgment-based evaluations remain narrower than general scientific value.

MLE-Ideator provides a more direct execution signal: an ideation model learns strategic guidance through reinforcement learning while a separate agent implements suggestions~\cite{mle-ideator}. Held-out MLE tasks test whether that guidance transfers, but objectives remain supplied. TTT-Discover updates model weights during search on a single problem, whereas Evolution Fine-Tuning trains on evolutionary trajectories before evaluation on held-out optimization tasks~\cite{ttt-discover,evolution-fine-tuning}. This contrast separates task-local adaptation from accumulated cross-task training. EFT's supplied objectives and predominantly shared search scaffold still limit what can be inferred about unfamiliar research workflows or taste.

\begin{table}[tb]
\centering
\small
\caption{What changes through experience? Mechanism and transfer are separate questions; representative works may use several mechanisms.}
\label{tab:learning-mechanisms}
\begin{tabularx}{\linewidth}{@{}P{0.18\linewidth}YY@{}}
\toprule
Mechanism & Changed object and temporal scope & Representative works \\
\midrule
Local reflection & Context and action selection during an investigation & Reflexion; Dolphin; OR-Agent \\
Memory and skills & External evidence or procedures reused in later sessions & AgentRxiv; Agora; EvoScientist; InternAgent-1.5; AutoResearchClaw; HypoForge \\
Program evolution & Controller code or research-loop strategy & ADAS; Darwin G\"odel Machine; Bilevel Autoresearch \\
Parameter training & Model weights, before or during a task & CycleResearcher; ResearchTrails; SciSense; MLE-Ideator; TTT-Discover; EFT \\
\bottomrule
\end{tabularx}
\end{table}

\subsection{Limits of current evidence and open questions}

Across these mechanisms, the central unresolved issue is attribution. A stronger final artifact may reflect additional attempts, retrieved solutions, a better implementer, or a better decision policy. Evaluation should therefore state what is retained, freeze the execution interface where possible, and compare both equal online budgets and total costs including experience collection. Resetting memory or weights between tasks must be explicit. Splitting individual trajectories is insufficient when training and test records share datasets, repositories, or near-identical task templates.

Three questions follow. Which experience preserves reasons and applicability conditions rather than copying successful answers? Can a researcher revise a previously useful strategy when the environment changes? Can learning improve the selection of informative experiments and stopping decisions, even when the final benchmark score does not improve? Answering these requires executed outcomes and independent checks, not only a judge that also supplies the training reward. The following pilots are proposed designs, not reported results; their task counts and budgets are deliberately modest starting assumptions.

\subsection{Proposed pilot A: the causal value of research memory}

\paragraph{Hypothesis and method.}
Structured decision records combining successful and unsuccessful attempts improve subsequent experimental choices more reliably than polished prose summaries. Create a fixed archive with states, proposed actions, stated reasons, observed outcomes, and unresolved uncertainties. Distinguish implementation errors, inconclusive measurements, and evidence against a hypothesis; do not collapse all three into failure. Render the same archive as prose or structured records, retaining provenance links. Freeze the archive before target evaluation so treatments cannot contaminate one another.

\paragraph{Data, splits, and controls.}
As a pilot design, select twelve small public-data ML tasks from tabular prediction, image classification, and compact language-model training. Allocate two tasks per family to archive construction, one to development, and one to a hidden target set; keep underlying datasets and data splits disjoint across these roles. Pre-register task identities after an execution smoke test. Randomize six conditions on the same target tasks and seeds: no memory, full prose summaries, success-only records, failure-only records, combined decision records, and task-shuffled combined records. The prose and combined-record comparison uses the same underlying evidence; success/failure comparisons isolate content. Match retrieval policy and memory-token limits across nonempty conditions. Shuffle task-to-memory associations rather than record order, preserving counts and lengths, to probe whether relevant experience matters.

\paragraph{Metrics and resource accounting.}
Use three seeds per target task. Record the independently checked test outcome of the final artifact selected on validation data, repeated-failure rate, redundant experiments, invalid conclusions, and time to a pre-specified validation improvement. A blinded audit of sampled decision logs checks whether selected experiments address the stated uncertainty. Cap each run at six experiments and 45 minutes of task computation on the same declared hardware: the main target comparison permits at most 40.5 device-hours if every run uses one device throughout. Two archive runs per source task and one development run per task and condition add ceilings of nine and 13.5 device-hours, respectively. Count summarization, retrieval, model tokens, and failed executions separately; also report actual utilization and API cost. These ceilings exclude remote model computation, which must remain in the cost ledger.

\paragraph{Failure risks.}
Task similarity could make retrieval an answer cache; outcome labels could incorrectly turn low power into a negative finding; long records could help through extra context alone. Shuffled-memory controls, provenance inspection, and matched evidence address parts of these risks. With only three target tasks, intervals should be grouped by task and treated as exploratory. A null effect would be informative if it separates useful failure prevention from improved research decisions.

\subsection{Proposed pilot B: distilling strategy across task families}

\paragraph{Hypothesis and method.}
Training on decision states, alternative actions, observed outcomes, and applicability conditions transfers better than imitating next actions alone. Use a fixed implementation agent and evaluator throughout; only the strategist changes. Start with a small open model and parameter-efficient supervised fine-tuning, postponing reinforcement learning until the supervision itself shows a benefit. Require every training outcome to link to an executed log. Unexecuted synthetic rationales may be retained as annotations, but never as observations or evidence that an action succeeded.

\paragraph{Data, splits, and comparisons.}
Propose sixteen tasks, four in each of four families: tabular prediction, image classification, compact language-model training, and numerical algorithm optimization. Reserve whole families, using two for training, one for development, and one for testing. Within every task, keep model-fitting data separate from validation feedback and a sealed final evaluator. The initial pilot uses one pre-registered family split; further rotations would be needed for a general transfer claim. Collect up to six short trajectories per training task, with four experimental decisions each, giving at most 192 executed decision transitions. Compare the base strategist, base plus retrieval from the same archive, action-only fine-tuning, and full decision-context fine-tuning. Use identical source transitions and optimizer-step budgets, report token counts, and add a token-matched development comparison; remove outcomes or applicability conditions in further development ablations.

\paragraph{Metrics, cost, and interpretation.}
Freeze all learned components on the held-out family. Evaluate four target tasks with three seeds, at most four experiments and 30 minutes per run, giving a 24 device-hour ceiling for the four-way target comparison under single-device execution. Set an initial fine-tuning cap of eight device-hours on declared hardware; report teacher trajectories, filtering, development, training, and inference costs separately, alongside performance after amortizing those costs over different numbers of future tasks. Primary outcomes are independently checked task improvement and improvement per total budget; secondary outcomes include invalid experiments, premature stopping, and sensitivity to a changed evaluator interface. Few training examples, family-specific conventions, and contamination can all mimic or hide transfer. Even a positive result would support reusable experimentation strategy under supplied objectives, not autonomous selection of important research questions.
